\documentclass[12pt]{article}
\usepackage[T2A]{fontenc}
\usepackage[utf8]{inputenc}
\usepackage[russian]{babel}
\usepackage{amsmath, amssymb, amsfonts}
\usepackage{geometry}
\usepackage{fancyhdr}
\usepackage{microtype}
\usepackage{titlesec} % Пакет для управления стилем заголовков
\usepackage{enumitem}
\usepackage{tikz}
\usepackage{tikz-3dplot}
\usetikzlibrary{calc, intersections, through, backgrounds, arrows.meta}

% Настройка Section: по центру, жирный, крупный
\pagestyle{fancy}
% Настройка полей
\geometry{a5paper, margin=1cm, top=2cm, headheight=15pt}

% --- ЖЕСТКАЯ КОРРЕКЦИЯ ЦЕНТРИРОВАНИЯ ---
\pagestyle{fancy}
\fancyhf{} % Полная очистка всех полей (убирает невидимые отступы)
\fancyhead[C]{\thepage} % Номер строго по центру
\renewcommand{\headrulewidth}{0pt} % Убираем линию
\fancyfoot{} % Низ пустой

% Убираем внутренние отступы колонтитулов, которые могут давать перекос
\setlength{\headwidth}{\textwidth} 

% Переопределяем системный стиль plain, чтобы он не сбивал настройки
\makeatletter
\let\ps@plain\ps@fancy
\makeatother

\parindent=0pt
\parskip=0.6em
\titleformat{\section}
  {\normalfont\bfseries\centering}{}{0pt}{}

% Настройка Subsection: отступ 1см, жирный, обычный размер
% {команда}[формат]{стиль}{метка}{отступ_от_метки}{перед_заголовком}[после_заголовка]
\titleformat{\subsection}
  {\normalfont\bfseries}{\hspace{1cm}}{1pt}{}
\titlespacing*{\subsection}{1cm}{2ex plus 1ex minus .2ex}{0.5ex plus .2ex}

% Убираем автоматическую нумерацию, чтобы заголовки были как в оригинале
\setcounter{secnumdepth}{0}
\AddEnumerateCounter{\asbuk}{\russian@asbuk@alph}{щ} % section 8.1
\setlist{wide, nosep, font=\rmfamily\normalfont}
\setlist[enumerate, 2]{label=\asbuk*)}
\setlist[enumerate, 3]{label=\asbuk*)}


\begin{document}

	\section{3.1. Геометрия Евклида}

	В этой главе даются задачи, касающиеся трех геометрий --- Евклида, Лобачевского и Римана. В каждой из этих геометрий(в соответсвующих моделях) описываются расстояния между точками, углы между прямыми, движения, расстояния между точками. В каждой из них изучается геометрия треугольника и тригонометрия.

	Пункт 3.1 посвящен теме, которую, в принципе, можно считать исчерпанной школьным образованием. Поэтому мы лишь кратко изложим темы, которые будут обсуждаться на примере неевклидовых геометрий.

	\begin{enumerate}
		\item Доказать, что любое движение евклидовой плоскости является композицией не более чем трех осевых симметрий.
		\item Доказать, что евклидовы движения, отражения и повороты, как суперпозиции параллельных переносов, отражений и поворотов.
		\item Доказать основные теоремы тригонометрии треугольника:

		Пусть $a$ и $b$ -- катеты, а $c$ -- гипотензуа прямоугольного треугольника $ABC$ с острыми углами $A$ и $B$. Тогда
		\begin{enumerate}
			\item $c^2 = a^2 + b^2$,
			\item $\tg A = \frac{a}{b}$,
			\item $\cos A = \frac{b}{c}$,
			\item $\sin A = \frac{a}{c}$,
			\item $\tg B = \frac{b}{a}$,
			\item $\cos B = \frac{a}{c}$,
			\item $\sin B = \frac{b}{c}$.
		\end{enumerate}

		Пусть в треугольнике стороны $a$, $b$, $c$ и углы $A$, $B$ и $C$. Тогда:
		\begin{enumerate}
			\item $\sin B = \frac{b}{c}$,
			\item $\cos B = \frac{a}{c}$.
		\end{enumerate}
		\item Любой невырожденный треугольник конгруэнтен треугольнику с координатами $(0,0)$, $(a,0)$, $(b\cos\varphi, b\sin\varphi)$, где $a>0$, $b>0$, $0<\varphi<\pi$.

		Найти его 1) медианы и их длины, 2) высоты и их длины, 3) биссектрисы и их длины, 4) центр описанного круга, 5) центр вписанного круга, 6) центр тяжести, 7) центр вневписанных окружностей.
		\newpage
		\item Доказать следующее классические теоремы:

		1) Теорему Герона; 2) Теорему Чевы; 3) Теорему Менелая.
	\end{enumerate}
	\section{3.2. Геометрии Лобачевского}

	\section{3.2.1. Дробно-линейные преобразования}

	Отображение вида $z \mapsto \frac{az + b}{cz + d}$, где $ad - bc \ne 0$, называется дробно-линейным преобразованием (область определения -- $\mathbb{C} \cup \{\infty\}$).

	\begin{enumerate}
		\item Доказать, что преобразование, обратное дробно-линейному, тоже дробно-линейно.
		\item Доказать, что дробно-линейные преобразования сохраняют углы.
		\item Доказать, что при добрно-линейном преобразовании прямая или окружность переходит в прямую или окружность.
		\item Доказать, что дробно-линейные преобразования сохраняют углы.
		\item Доказать, что если точки $z_1, z_2, z_3$ попарно различны и точки $w_1, w_2, w_3$ попарно различны, то существует дробно-линейное преобразование, переводящее $z_i$ в $w_i$.
		\item Доказать, что при дробно-линейном преобразовании сохраняется двойное отношение четырех точек $ \frac{z_1 - z_3}{z_1 - z_3} : \frac{z_2 - z_3}{z_2 - z_4}$.
		\item
		\begin{enumerate}
			\item Доказать, что точки $z_1, z_2, z_3$ ляжат на одной прямой тогда и только тогда, когда $\frac{z_1 - z_2}{z_1 - z_3} \in \mathbb{R}$.
			\item Доказать, что точки $z_1, z_2, z_3, z_4$ лежат на одной прямой или окружности тогда и только тогда, когда их двойное отношение вещественно.
		\end{enumerate}
		\item Приведите пример дробно-линейного преобразования, переводящегоо верхнюю поуплоскость $\operatorname{Im} z \ge 0$ в единичный круг $|z| \le 1$.
		\item
		\begin{enumerate}
			\item Доказать, что дробно-линейное преобразование полностью задается образами трех различных точек.
			\item Доказать, что если дробно-линейное преобразование оставляет неподвижными три различные точки, то оно тождественно
		\end{enumerate}
		\item Доказать, что преобразование, сохраняющее двойное отношение любых четырех точек, является дробно-линейным.
	\end{enumerate}
	\section{3.2.2. Преобразования верхней полуплоскости}
	\begin{enumerate}[resume*]
		\item Пусть $a, b, c, d \in \mathbb{R}$. Доказать, что преобразования $z \mapsto \frac{az + b}{cz + d}$, где $ad - bc > 0$ и $z \mapsto \frac{a\overline z + b}{c\overline z + d}$, где $ad - bc < 0$, сохраняет верхнюю полуплоскость $\operatorname{Im} z > 0$.
	\end{enumerate}
	Обозначим множество вещественных матриц порядка 2 с определителем 1 через $\operatorname{SL}(2, \mathbb{R})$.
	\begin{enumerate}[resume*]
		\item Доказать, что если $\operatorname{Im} z_1 > 0$ и $\operatorname{Im} z_2 > 0$, то существует преобразование вида $z \mapsto \frac{az + b}{cz + d}$, где $\begin{pmatrix} a & b \\ c & d \end{pmatrix} \in \operatorname{SL}(2, \mathbb{R})$, переводящее $z_1$ в $z_2$.

		\item Преобразование $z \mapsto \frac{az + b}{cz + d}$, где $\begin{pmatrix} a & b \\ c & d \end{pmatrix} \in \operatorname{SL}(2, \mathbb{R})$, оставляет неподвижными точки $z_1$ и $z_2$, лежащие в верхней полуплоскости. Доказать, что оно тождественно.

		\item Рассмотрим все преобразования $z \mapsto \frac{az + b}{cz + d}$, где $\begin{pmatrix} a & b \\ c & d \end{pmatrix} \in \operatorname{SL}(2, \mathbb{R})$ сохраняющие точку $i$. Доказать, что образы фиксированной точки $z_0$ при таких преобразованиях лежат на одной окружности.
	\end{enumerate}
	\section{3.2.3. Расстояние на плоскости Лобачевского}

	\textit{Определение}. Для модели Пуанкаре геометрии Лобачевского в верхней полуплоскости расстояние между точками $(x_1, x_2)$ и $(y_1, y_2)$ равно $\operatorname{arch} \left(1 + \frac{(x_1 - y_1)^2 + (x_2 - y_2)^2}{2 x_2 y_2}\right)$.
	\begin{enumerate}[resume*]
		\item Доказать, что расстояние сохраняется при преобразованиях указанных в задаче 11.
		\item
		\begin{enumerate}
			\item Доказать, что для точек $(0, x_2)$ и $(0, y_2)$ расстояние равно $\ln |x_2 / y_2| = |\ln[0, \infty, x_2, y_2 ]|$, где $[0, \infty, x_2, y_2]$ -- двойное отношение указанных точек.

			\item Окружность с центром на оси $Ox$, проходящая через точки $z$ и $w$, пересекает ось $Ox$ в точках $A$ и $B$. Доказать, что расстояние между точками $z$ и $w$ равно $|\ln [A, B, z, w]|$.
		\end{enumerate}
		\item Доказать, что если точка $C$ лежит на отрезке $AB$, то $AB = AC + BC$.

		\item Доказать, что минимальное расстояние от точки $(x_1, x_2)$ до точки луча $Oy$ равно $\operatorname{arch} \frac{\sqrt{x_1^2 + x_2^2}}{x_2}$; это расстояние равно длине перпендикуляра, опущенного из точки на прямую.

		\item 
		\begin{enumerate}
			\item Доказать, что при движении точки прямой $x=1$ по направлению к оси $Ox$ расстояние от нее до прямой $Oy$ бесконечно возрастает.

			\item Доказать, что при движении точки вдоль (гиперболической) прямой $(x-1)^2 + y^2 = 1$ по направлению к точке $(0,0)$ расстояние от нее до прямой $Oy$ стремится к нулю.

		\end{enumerate}
		\item Доказать, что гиперболическая окружность в модели Пуанкаре является евклидовой окружностью (с другим центром). Если центр гиперболической окружности лежит на оси $Oy$, то центр евклидовой окружности тоже лежит на оси $Oy$.
		\item Доказать, что $AB + BC \ge AC$, причем равенство достигается лишь в том случае, когда точка $B$ лежит на отрезке $AC$.

		\item Доказать, что все точки евклидовой прямой $y = kx$ равнодуальны от гиперболической прямой $Oy$.

		\item
		\begin{enumerate}
			\item Доказать, что множество точек, равноудаленных от двух данных точек, представляет собой прямую.

			\item Доказать, что любое движение плоскости Лобачевского можно представить в виде композиции не более чем трех симметрий относительно прямых.
		\end{enumerate}
	\end{enumerate}
	\textit{Определение}. Прямые, имеющие общую точку на оси $Ox$ или общую бесконечно удаленную точку, называются параллельными. Не пересекающиеся и не параллельные прямые называются расходящимися.
	\begin{enumerate}[resume*]
		\item 
		\begin{enumerate}
			\item Доказать, что любые две расходящиеся прямые имеют единственный общий перпендикуляр.

			\item Доказать, что если две прямые имеют общий перпендикуляр, то они расходятся.
		\end{enumerate}
		\item Доказать, что множество точек, равноудаленных от данных прямых, представляет собой одну или две прямые.
		\newpage
		\item Доказать, что любую пару параллельных прямых можно перевести в любую другую пару параллельных прямых.
	\end{enumerate}
	\section{3.2.4. Углы в геометрии Лобачевского}
	
	Угол между соответствующими прямыми или полуплоскостями в верхней полуплоскости (в верхней полуплоскости или в единичном круге), согласно задаче 4 движения Лобачевского углы.
	\begin{enumerate}[resume*]
		\item Доказать, что сумма углов треугольника меньше $\pi$.

		\item Треугольник $T$ разрезан на треугольники $T_1$ и $T_2$. Доказать, что $\sigma(T) = \sigma(T_1) + \sigma(T_2)$, гдк $\sigma(T) = \pi - \text{сумма углов $T$}$.

		\item Доказать, что если углы двух треугольников равны, то равны сами треугольники.

		\item Доказать, что площадь треугольника с углами 
		\begin{enumerate}
			\item $0, 0, 0$ равна $\pi$;
			\item $0, \alpha, \beta$ равна $\pi = (\alpha + \beta)$;
			\item$\alpha, \beta, \gamma$ равна $\pi - (\alpha + \beta + \gamma)$.
		\end{enumerate}
	\end{enumerate}
	\section{3.2.5 Ортогональные проекции}
	\begin{enumerate}[resume*]
		\item
		\begin{enumerate}
			\item Доказать, что ортогональная проекция одной из сторон острого угла на другую сторону представляет собой ограниченное множество.

			\item доказать, что ортогональная проекция одной из расходящихся прямых на другую представляет собой ограниченное множество.
		\end{enumerate}
	\end{enumerate}
	\section{3.2.6. Длина окружности}
	\begin{enumerate}[resume*]
		\item Доказать, что при отображении $z \mapsto \frac{z-i}{z+i}$ метрика $\frac{dx^2 + dy^2}{y^2}$ переходит в $\frac{4|dz^2|}{(1-|z|^2)^2}$.

		\item Доказать, что в модели Пуанкаре в единичном круге окружность радиуса $r$ с центром в центре круга представляется собой евклидову окружность радиуса $R = \frac{e^r - 1}{e^r + 1} = \th\left( \frac{r}{2} \right)$.

		\item Доказать, что длина окружности радиуса $r$ равна $2 \pi \sh r$.
	\end{enumerate}
	\section{3.2.7. Гиперболическая тригонометрия}

	Определение. Пусть $AH$ --- перпендикуляр к прямой $l$, а луч $AB$ параллелен $l$. Тогда угол $\alpha = \angle HAB$ зависит лишь от длины $a$ отрезка $AH$; он называется углом параллельности (в обозначениях Лобачевского $\alpha = \Pi(a)$).

	\begin{enumerate}[resume*]
		\item Доказать, что $a$ и $\alpha$ связаны следующими соотношениями:
		\begin{enumerate}
			\item $\ch a \sin \alpha = 1$;
			\item $\sh a \tg \alpha = 1$;
			\item $\th a = \cos \alpha$;
			\item $e^{-a} = \tg \frac{\alpha}{2}$.
		\end{enumerate}

		\item Доказать, что стороны и углы прямоугольного треугольника с прямым углом $\gamma$ связаны следующими соотношениями:
		\begin{enumerate}
			\item $\ch c = \ch a \ch b$;
			\item $\th a = \sh b \tg \alpha$;
			\item $\ctg \alpha \ctg \beta = \ch c$;
			\item $\sh a = \sh c \sin \alpha$;
			\item $\th b = \th c \cos \alpha$;
			\item $\cos \alpha = \ch a \sin \beta$.
		\end{enumerate}
	\end{enumerate}

	\textit{Замечание}. Формулы а), б), г), и д) имеют аналоги в евклидовой тригонометрии (см. задачу 3 \S 3.1). Формулы в) и е) специфичны для геометрии Лобачевского, ибо сумма углов $\alpha + \beta$ не равна константе.
	\begin{enumerate}[resume*]
		\item Доказать, что в треугольнике противолежащими им углами $\alpha, \beta, \gamma$ выполняются следующие соотношения:
		\begin{enumerate}
			\item $\frac{\sh a}{ \sin \alpha} = \frac{\sh b}{ \sin beta} = \frac{\sh c}{ \sin \gamma}$ (теорема синусов);
			\item $\ch a = \ch b \ch c - \sh b \sh c \cos \alpha$ (первая теорема косинусов);
			\item $\ch a \sin \beta = \ch b \sin \alpha \cos \gamma + \cos \alpha \sin \gamma$;
			\item $\ch a = \frac{\cos \alpha + \cos \beta \cos \gamma}{\sin \beta \sin \gamma}$ (вторая теорема косинусов).
		\end{enumerate}

		\item С помощью задачи 37, б) доказать, что $a < b + c$

		\item С помощью задачи 37, г) доказать, что если соотвественные углы двух треугольников равны, то равны и их соответственные стороны.

		\item 
		\begin{enumerate}
			\item Доказать, что при $\alpha < \frac{\pi(n-2)}{n}$ существует правильный $n$-угольник, все углы которого равны $\alpha$ и все стороны которого тоже равны.
			\item Доказать, что если $\alpha$ и $r$ -- сторона и радиус описанной окружности правильного $n$-угольника с углом $\alpha$, то
				\[
					\ch r = \frac{\cos \frac{\alpha}{2}}{\sin \frac{\pi}{n}} \text{ и } \ch \frac{a}{2} = \frac{\cos \frac{\pi}{n}}{\sin \frac{\alpha}{2} }
				\]
				.
			\item При каких $n$ существует правильный $n$-угольник с углом 

			$\alpha = \frac{2\pi}{n}$(т.е. суммой углов $2 \pi$)?
		\end{enumerate}

		\item Три угла четырехугольника равны $\frac{\pi}{2}$, а четвертый угол равен $\gamma$. Найти $\gamma$, если известны длины $a$ и $b$ сторон, соединяющих прямые углы.
	\end{enumerate}


	\section{3.2.8. Геометрия треугольника}
 
	\begin{enumerate}[resume*]
		\item (Теорема Менелая) Прямая $l$ пересекает стороны $BC$, $CA$ и $AB$ треугольника $ABC$ (или их продолжения) в точках $A_1$, $B_1$ и $C_1$. Доказать, что
		\[
			\frac{\sh AC_1}{\sh C_1 B} \cdot \frac{\sh BA_1}{\sh A_1 C} \cdot \frac{\sh CB_1}{\sh B_1 A} = 1.
		\]
 		\begin{center}
		\begin{tikzpicture}[scale=1.0, line cap=round, line join=round]
				% Основные точки
			  \coordinate (A)  at (0,0);
			  \coordinate (B)  at (1.35,3.05);
			  \coordinate (C)  at (4.20,0.42);
			  \coordinate (B1) at (4.92,0.46); % B_1 на продолжении AC

			  % Точки на гиперболической линии l (без перегибов)
			  \coordinate (A2) at (0.00,1.72);
			  \coordinate (A2e) at (0.00,1.90);
			  \coordinate (C2) at (0.95,2.10);
			  \coordinate (A1) at (1.97,2.20);
			  \coordinate (C1) at (4.20,1.15);
			  \coordinate (C1e) at (4.20,1.26);
			  \coordinate (B2) at (1.35,1.64);

			  % "Треугольник" (кривые стороны)
			  \draw[thick] (A) .. controls (0.62,1.12) and (1.00,2.06) .. (B);
			  \draw[thick] (B) .. controls (2.00,1.78) and (3.05,1.00) .. (C);
			  % Нижняя сторона и ее продолжение: одна кривая без стыка в C
			  \draw[thick,domain=-0.18:4.92,samples=150,smooth]
				plot (\x,{-0.00903342*(\x)^2 + 0.13794*(\x)});

			  % Линия l: одна коническая (параболическая) дуга без точек перегиба
			  \draw[thick,domain=-0.18:4.92,samples=130,smooth]
				plot (\x,{-0.167934*(\x)^2 + 0.570137*(\x) + 1.72});

			  % Вспомогательные отрезки
			  \draw[thick] (A) -- (A2);
			  \draw[thick] (C) -- (C1e);
			  \draw[thick] (B) -- (B2); % BB2

			  % Подписи (с вынесением от кривых)
			  \node[left,xshift=-3pt,yshift=-1pt]  at (A) {$A$};
			  \node[above,yshift=4pt]              at (B) {$B$};
			  \node[below,yshift=-4pt]             at (C) {$C$};

			  \node[left,xshift=-4pt,yshift=2pt]   at (A2) {$A_2$};
			  \node[left,xshift=-4pt,yshift=4pt]   at (C2) {$C_1$};
			  \node[above,yshift=5pt]              at (A1) {$A_1$};
			  \node[right,xshift=5pt,yshift=2pt]   at (C1) {$C_2$};
			  \node[right,xshift=6pt,yshift=1pt]   at (B1) {$B_1$};
			  \node[right,xshift=5pt,yshift=-1pt]  at (B2) {$B_2$};
			  \node[above,xshift=0pt,yshift=4pt]   at (3.05,1.68) {$l$};
		\end{tikzpicture}
		\end{center}
		\item (Теорема Чевы) На сторонах $BC$, $CA$ и $AB$ треугольника $ABC$ (но не на их продолжениях!) взяты точки $A_1$, $B_1$ и $C_1$. Доказать, что отрезки $AA_1$, $BB_1$ и $CC_1$ пересекаются в одной точке тогда и только тогда, когда выполнено одно из следующих эквивалентных условий:
		\begin{enumerate}
			\item $\dfrac{\sin ACC_1}{\sin C_1CB} \cdot \dfrac{\sin BAA_1}{\sin A_1AC} \cdot \dfrac{\sin CBB_1}{\sin B_1BA} = 1$;
			\item $\dfrac{\sh AC_1}{\sh C_1 B} \cdot \dfrac{\sh BA_1}{\sh A_1 C} \cdot \dfrac{\sh CB_1}{\sh B_1 A} = 1$.
		\end{enumerate}
 
		\item Доказать, что в одной точке пересекаются:
		\begin{enumerate}
			\item биссектрисы любого треугольника;
			\item медианы любого треугольника;
			\item высоты остроугольного треугольника.
		\end{enumerate}
	\end{enumerate}
 
	\textit{Замечание}. В гиперболической геометрии есть три типа пучков прямых: 1) прямые, пересекающиеся в одной точке; 2) параллельные прямые; 3) прямые, перпендикулярные одной прямой. В каждом из этих случаев будем говорить, что прямые принадлежат одному пучку.
 
	\begin{enumerate}[resume*]
		\item Доказать, что перпендикуляры к сторонам треугольника, прохходящие через их середины, принадлежат одному пучку.
 
		\item
		\begin{enumerate}
			\item Доказать, что высоты треугольника принадлежат одному пучку.
			\item Доказать, что биссектрисы двух внешних углов и одного внутреннего угла треугольника принадлежат одному пучку.
		\end{enumerate}
 
		\item Доказать формулу Герона:
		\[
			\sin\left(\frac{\Delta}{2}\right) = \frac{\sqrt{(e^s - 1)(e^{s-2a} - 1)(e^{s-2b} - 1)(e^{s-2c} - 1)}}{(e^{a} + 1)(e^{b} + 1)(e^{c} + 1)},
		\]
		где $\Delta$ --- площадь треугольника, $s = a + b + c$.
	\end{enumerate}
 
	\section{3.2.9. Модель Клейна}
 
	Модель Клейна получается из модели Пуанкаре в единичном круге следующим образом. Представим единичный круг как экватор сферы и спроецируем экваториальную плоскость из северного полюса, а затем точки южного полушария ортогонально спроецируем на экваториальную полуплоскость. В результате окружности, перпендикулярные границе единичного круга, переходят сначала в сечения сферы плоскостями, перпендикулярными экватору, а затем в хорды единичного круга.
 	
	Таким образом, дуге окружности, проходящей через точки $A$ и $B$ граничной окружности (и ортогональной ей), соответсвует хорда $AB$.
	\begin{enumerate}[resume*]
 			
		% \begin{center}
		% \begin{tikzpicture}[scale=1.3]
		% 	% Draw the sphere as an ellipse (oblique projection)
		% 	\draw[thick] (0,0) ellipse (1.8cm and 1.8cm);
		% 	% Equator as ellipse (slightly flattened to suggest 3D)
		% 	\draw[thick] (0,0) ellipse (1.8cm and 0.55cm);
		% 	% North pole S
		% 	\coordinate (S) at (0, 1.8);
		% 	\filldraw (S) circle (1.5pt) node[above right] {$S$};
		% 
		% 	% Points A and B on the equator (inside the disk view)
		% 	\coordinate (Apt) at (-0.7, -0.12);  % A on equator ellipse
		% 	\coordinate (Bpt) at (0.5, -0.08);   % B on equator ellipse
		% 	\filldraw (Apt) circle (1pt) node[left] {$A$};
		% 	\filldraw (Bpt) circle (1pt) node[right] {$B$};
		% 
		% 	% Chord AB
		% 	\draw[thick] (Apt) -- (Bpt);
		% 
		% 	% Points Z and W on the sphere surface (above equator)
		% 	\coordinate (Z) at (-0.35, 0.65);
		% 	\coordinate (W) at (0.55, 0.8);
		% 	\filldraw (Z) circle (1pt) node[left] {$Z$};
		% 	\filldraw (W) circle (1pt) node[above right] {$W$};
		% 
		% 	% Dashed lines from S through Z and W to the equatorial disk
		% 	\draw[dashed] (S) -- (Apt);
		% 	\draw[dashed] (S) -- (Bpt);
		% 
		% 	% Lines from S through Z and W showing projection
		% 	\draw[thin] (S) -- ($(S)!1.55!(Z)$);
		% 	\draw[thin] (S) -- ($(S)!1.45!(W)$);
		% 
		% 	% Projected points Z', W' on equatorial plane (below Z, W)
		% 	\coordinate (Zp) at ($(S)!1.55!(Z)$);
		% 	\coordinate (Wp) at ($(S)!1.45!(W)$);
		% 	\filldraw (Zp) circle (1pt) node[below] {$Z'$};
		% 	\filldraw (Wp) circle (1pt) node[right] {$W'$};
		% \end{tikzpicture}
		% \end{center}

		\item В модели Пуанкаре прямая, проходящая через точки $Z$ и $W$, пересекает граничную окружность в точках $A$ и $B$. В модели клейна точкам $Z$ и $W$ соотвествуют точки $Z'$ и $W$.

		Доказать, что $\ln[A, B, Z', W'] = 2\ln[A, B, Z, W].$
 
		\item Доказать, что точке $z$ в модели Пуанкаре соответствует точка $w = \frac{2z}{|z|^2 + 1}$ в модели Клейна.


		\item Доказать, что в модели Клейна прямые принадлежат одному пучку тогда и только тогда, когда содержащие их евклидовы пряме имеют общую точку.
 
		\item Доказать, что изображённые на рис. а) и б) построения позволяют построить середину $M$ отрезка $AB$ и биссектрису угла между прямыми $a$ и $b$.

		\begin{center}
		\begin{tikzpicture}[scale=1.25]
			%--- Figure a): midpoint construction in Klein disk ---
			\begin{scope}
				% Unit disk boundary
				\draw[thick] (0,0) circle (1.5cm);
				% Points A and B on the chord
				\coordinate (A) at (-0.9, -0.5);
				\coordinate (B) at (0.85, -0.3);
				\filldraw (A) circle (1pt) node[left] {$A$};
				\filldraw (B) circle (1pt) node[right] {$B$};
				% The chord AB (hyperbolic line)
				% Find where line AB intersects the boundary circle radius 1.5
				% Direction vector
				% We extend AB to boundary: parametric A + t*(B-A)
				% |A + t*(B-A)|^2 = 1.5^2  => solve for t
				% A=(-0.9,-0.5), B=(0.85,-0.3), B-A=(1.75,0.2)
				% (-0.9+1.75t)^2 + (-0.5+0.2t)^2 = 2.25
				% 0.81 - 3.15t + 3.0625t^2 + 0.25 - 0.2t + 0.04t^2 = 2.25
				% 3.1025t^2 - 3.35t - 1.19 = 0  => t=(3.35 ± sqrt(11.2225+14.7628))/6.205
				% t=(3.35 ± 5.091)/6.205 => t1=1.362, t2=-0.280
				\coordinate (P1) at ($(A)!1.362!(B)$);   % boundary point beyond B
				\coordinate (P2) at ($(A)!-0.280!(B)$);  % boundary point beyond A
				\draw[thick] (P2) -- (P1);
 
				% Four lines from A and B to opposite boundary points, forming an X
				% Lines through A: from P1 through A to boundary, and from P2 through B to boundary
				% We draw the diagonals of the complete quadrilateral to find M
				% Two extra points on boundary top area
				\coordinate (T1) at (-0.6, 1.37);   % top-left on boundary
				\coordinate (T2) at (0.6,  1.37);   % top-right on boundary
 
				% Lines: T1--P1 (passes through somewhere) and T2--P2
				% The intersection of diagonals gives M
				% Line T1-A extended, T2-B extended
				\draw[thin] (T1) -- (P1);
				\draw[thin] (T2) -- (P2);
				% Cross lines: T1-B and T2-A to boundary
				\draw[thin] (T2) -- (P2);
				% Actually draw the 4 lines of the complete quadrilateral:
				% Through A: T1-A and T2-A
				\draw[thin] (T1) -- ($(T1)!2.2!(A)$);
				\draw[thin] (T2) -- ($(T2)!2.1!(A)$);
				% Through B: T1-B and T2-B  
				\draw[thin] (T1) -- ($(T1)!2.0!(B)$);
				\draw[thin] (T2) -- ($(T2)!1.9!(B)$);
 
				% Midpoint M = intersection of diagonals T1B∩T2A  (approx)
				% T1=(-0.6,1.37), B=(0.85,-0.3): dir=(1.45,-1.67), param: (-0.6+1.45t, 1.37-1.67t)
				% T2=(0.6,1.37),  A=(-0.9,-0.5): dir=(-1.5,-1.87)
				% Intersection: -0.6+1.45t = 0.6-1.5s => 1.45t+1.5s=1.2
				%               1.37-1.67t = 1.37-1.87s => 1.67t=1.87s => s=1.67t/1.87=0.893t
				% 1.45t + 1.5*0.893t = 1.2 => 1.45t+1.339t=1.2 => 2.789t=1.2 => t=0.430
				\coordinate (M) at (-0.377, 0.652);  % approx
				\filldraw[fill=white] (M) circle (1.5pt) node[right] {$M$};
 
				% Tick marks on AM and MB on the chord
				\draw[thin] ($(A)!0.5!(M)$) ++(0.05,-0.08) -- +(-0.1,0.08);
				\draw[thin] ($(M)!0.5!(B)$) ++(0.05,-0.08) -- +(-0.1,0.08);
 
				\node[below] at (0,-1.7) {а)};
			\end{scope}
 
			%--- Figure b): angle bisector construction in Klein disk ---
			\begin{scope}[xshift=4cm]
				\draw[thick] (0,0) circle (1.5cm);
 
				% Two chords (hyperbolic lines a and b) intersecting inside the disk
				% Line a: from boundary point Qa1 to Qa2
				\coordinate (Qa1) at (-1.3, 0.7);
				\coordinate (Qa2) at (0.9, -1.2);
				% Line b: from boundary point Qb1 to Qb2
				\coordinate (Qb1) at (-1.0, -1.1);
				\coordinate (Qb2) at (1.2,  0.9);
 
				\draw[thick] (Qa1) -- (Qa2) node[below right, font=\small] {$a$};
				\draw[thick] (Qb1) -- (Qb2) node[above right, font=\small] {$b$};
 
				% Intersection point I of lines a and b
				% Line a: (-1.3,0.7)+t*(2.2,-1.9)
				% Line b: (-1.0,-1.1)+s*(2.2,2.0)
				% -1.3+2.2t = -1.0+2.2s => 2.2t-2.2s=0.3 => t-s=0.136
				%  0.7-1.9t = -1.1+2.0s => -1.9t-2.0s=-1.8 => 1.9t+2.0s=1.8
				% s=t-0.136 => 1.9t+2.0(t-0.136)=1.8 => 3.9t=2.072 => t=0.531
				\coordinate (I) at (-0.132, -0.310);
				\filldraw (I) circle (1pt);
 
				% Bisector construction: connect opposite boundary endpoints
				% Diagonal 1: Qa1 -- Qb2
				\draw[thin] (Qa1) -- (Qb2);
				% Diagonal 2: Qa2 -- Qb1
				\draw[thin] (Qa2) -- (Qb1);
 
				% The two bisectors pass through I and intersection of diagonals
				% Intersection of Qa1-Qb2 and Qa2-Qb1 (this gives the "pole"):
				% Qa1=(-1.3,0.7) dir to Qb2=(1.2,0.9): dir=(2.5,0.2), param (-1.3+2.5t,0.7+0.2t)
				% Qa2=(0.9,-1.2) dir to Qb1=(-1.0,-1.1): dir=(-1.9,0.1), param (0.9-1.9s,-1.2+0.1s)
				% -1.3+2.5t=0.9-1.9s => 2.5t+1.9s=2.2
				%  0.7+0.2t=-1.2+0.1s => 0.2t-0.1s=-1.9 => s=2t+19  (approximate issue, recompute)
				% 0.7+0.2t = -1.2+0.1s => 0.1s=1.9+0.2t => s=19+2t
				% 2.5t+1.9(19+2t)=2.2 => 2.5t+36.1+3.8t=2.2 => 6.3t=-33.9 => off, lines nearly parallel
				% Use other diagonal pair for bisectors:
				% Bisector 1: line through I and the intersection of Qa1-Qb1 extended to boundary
				% For simplicity, show approximate bisector as line through I toward boundary:
				\coordinate (Bis1a) at (-1.45, -0.58);
				\coordinate (Bis1b) at (1.2, 0.3);
				\draw[dashed, thin] (Bis1a) -- (Bis1b);
 
				\coordinate (Bis2a) at (0.05, -1.5);
				\coordinate (Bis2b) at (-0.3, 1.49);
				\draw[dashed, thin] (Bis2a) -- (Bis2b);
 
				\node[below] at (0,-1.7) {б)};
			\end{scope}
		\end{tikzpicture}
		\end{center}
	\end{enumerate}
 
	\section{3.3. Геометрия Римана}
 
	\begin{enumerate}
		\item
		\begin{enumerate}
			\item Доказать, что множество точек сферы, равноудалённых от двух данных точек сферы, представляет собой прямую (окружность большого круга).
 
			\item Доказать, что любое движение сферы можно представить в виде композиции не более чем трёх симметрий относительно прямых.
		\end{enumerate}
 
		\item Найти площадь сферического треугольника с углами $A$, $B$ и $C$, если радиус сферы равен $R$.
 
		\item
		\begin{enumerate}
			\item На сфере фиксированы точки $A$ и $B$. Найти множество точек $C$, для которых $\angle A + \angle B - \angle C = \mathrm{const}$.
 
			\item На сфере фиксированы точки $A$ и $B$. Найти множество точек $C$, для которых $S_{ABC} = \mathrm{const}$.
		\end{enumerate}
 
		\textit{Обозначение}: $a = BC$ -- угловая величина стороны $BC$, $A$ -- величина угла при вершине $A$.
 
		\item
		\begin{enumerate}
			\item Доказать, что $a < b + c$.
			\item Доказать, что $a + b + c < 2\pi$.
			\item Доказать, что $A + B + C > \pi$.
		\end{enumerate}
 
		\item Доказать, что
		\begin{enumerate}
			\item $\dfrac{\sin A}{\sin a} = \dfrac{\sin B}{\sin b} = \dfrac{\sin C}{\sin c}$ (теорема синусов)
			\item $\cos a = \cos b \cos c + \sin b \sin c \cos A$ (1-ая теорема косинусов)
			\item $\cos A = -\cos B \cos C + \sin B \sin C \cos a$ (2-ая теорема косинусов)
		\end{enumerate}
 
		\item
		\begin{enumerate}
			\item Доказать, что если $a, b, c > \frac{\pi}{2}$, то $A, B, C > \frac{\pi}{2}$.
			\item Доказать, что если $A, B, C < \frac{\pi}{2}$, то $a, b, c < \frac{\pi}{2}$.
		\end{enumerate}
 
		\item Доказать, что если $a, b, c < \frac{\pi}{2}$, то треугольник $ABC$ можно разместить внутри треугольника, все углы которого прямые.
 
		\item (Теорема Менелая). Прямая пересекает стороны треугольника $ABC$ (или их продолжения) в точках $A_1$, $B_1$ и $C_1$. Доказать, что
		\[
			\frac{\sin AC_1}{\sin C_1 B} \cdot \frac{\sin BA_1}{\sin A_1 C} \cdot \frac{\sin CB_1}{\sin B_1 A} = 1.
		\]
 
		\item (Теорема Чевы) На сторонах треугольника $ABC$ взяты точки $A_1$, $B_1$ и $C_1$. Доказать, что отрезки $AA_1$, $BB_1$ и $CC_1$ пересекаются в одной точке тогда и только тогда, когда
		\[
			\frac{\sin AC_1}{\sin C_1 B} \cdot \frac{\sin BA_1}{\sin A_1 C} \cdot \frac{\sin CB_1}{\sin B_1 A} = 1.
		\]
	\end{enumerate}
 
	\section{3.4. Проективная геометрия}
 
	\begin{enumerate}
		\item Докажите, что существует проективное отображение, которое три данные точки одной прямой переводит в три данные точки другой прямой.
 
		\item Докажите, что проективное преобразование прямой однозначно определяется образами трёх произвольных точек.
 
		\item Докажите, что преобразование $P$ числовой прямой является проективным тогда и только тогда, когда оно представляется в виде
		\[
			P(x) = \frac{ax + b}{cx + d},
		\]
		где $a, b, c, d$ --- такие числа, что $ad - bc \ne 0$. (Такие отображения называют \textit{дробно-линейными}.)
 
		\textit{Определение}. Пусть $\alpha_1$ и $\alpha_2$ --- две плоскости в пространстве, $O$ --- точка, не лежащая ни на одной из этих плоскостей. \textit{Центральным проектированием} плоскости $\alpha_1$ на плоскость $\alpha_2$ с центром $O$ называют отображение, которое точке $A_1$ плоскости $\alpha_1$ ставит в соответствие точку пересечения прямой $OA_1$ с плоскостью $\alpha_2$.
 
		Прямую, на которой не определено центральное проектирование, называют \textit{исключительной прямой} данной проекции.
 
		\item Докажите, что при центральном проектировании прямая, не являющаяся исключительной, проецируется в прямую.
 
		Для того чтобы центральное проектирование было определено всюду, удобно считать, что на каждой прямой помимо обычных точек имеется ещё одна точка, называемая \textit{бесконечно удалённой}. При этом, если две прямые параллельны, то их бесконечно удалённые точки совпадают, другими словами, параллельные прямые пересекаются в бесконечно удалённой точке.
 
		Мы также будем считать, что на каждой плоскости помимо обычных прямых имеется ещё \textit{бесконечно удалённая прямая}, на которой лежат все бесконечно удалённые точки прямых данной плоскости. Бесконечно удалённая прямая пересекается с каждой обычной прямой $l$, лежащей на той же плоскости, в бесконечно удалённой точке прямой $l$.
 
		Если ввести бесконечно удалённые точки и прямые, то центральное проектирование плоскости $\alpha_1$ на плоскость $\alpha_2$ с центром в точке $O$ будет определено для всех точек плоскости $\alpha_1$. При этом исключительная прямая будет проецироваться в бесконечно удалённую прямую плоскости $\alpha_2$, а именно, образом точки $M$ исключительной прямой будет бесконечно удалённая точка прямой $OM$; в этой точке пересекаются прямые плоскости $\alpha_2$, параллельные прямой $OM$.
 
		\item Докажите, что если наряду с обычными точками и прямыми рассматривать бесконечно удалённые, то
		\begin{enumerate}
			\item через любые две точки проходит единственная прямая;
			\item любые две прямые, лежащие в одной плоскости, пересекаются в единственной точке;
			\item центральное проектирование одной плоскости на другую является взаимно однозначным отображением.
		\end{enumerate}
 
		\textit{Определение}. Отображение $P$ плоскости $\alpha$ на плоскость $\beta$ называется \textit{проективным}, если оно является композицией центральных проектирований и аффинных преобразований, т.е. если существуют плоскости $\alpha_0 = \alpha, \alpha_1, \ldots, \alpha_n = \beta$ и отображения $P_i$ плоскостей $\alpha_i$ на $\alpha_{i+1}$, каждое из которых является либо центральным проектированием, либо аффинным преобразованием, причём $P$ является композицией преобразований $P_i$. В случае, когда плоскость $\alpha$ совпадает с плоскостью $\beta$, отображение $P$ называют \textit{проективным преобразованием} плоскости $\alpha$. Прообраз бесконечно удалённой прямой мы будем называть \textit{исключительной} прямой данного проективного преобразования.
 
		\item
		\begin{enumerate}
			\item Докажите, что проективное преобразование $P$ плоскости, переводящее бесконечно удалённую прямую в бесконечно удалённую прямую, является аффинным.
 
			\item Докажите, что если точки $A$, $B$, $C$, $D$ лежат на прямой, параллельной исключительной прямой проективного преобразования $P$ плоскости $\alpha$, то $P(A)P(B) : P(C)P(D) = AB : CD$.
		\end{enumerate}
 
		\item
		\begin{enumerate}
			\item Докажите, что проективное преобразование $P$ переводит аффинные преобразования, либо это его исключительная прямая параллельна прямым $l_1$ и $l_2$, то $P$ переводит конечных и бесконечных точек каждую прямую прямую.
		\end{enumerate}
 
		\item Докажите, что если существует проективное преобразование, которое переводит лежащую внутри окружности, переводит в центр образа, а данную окружность --- в окружность, а исключительная прямая перпендикулярна данному диаметру, проходящему через точку $M$ --- в её центр, то это --- поворот или симметрия.
 
		\item Проективное преобразование некоторую окружность переводит в себя, а её центр остаётся на месте. Докажите, что это --- поворот или симметрия.
	\end{enumerate}
 
	\section{3.5. Ответы, указания и решения}
 
	\subsection{3.2.1. Дробно-линейные преобразования}
 
	\begin{enumerate}
		\item Указание. Прямые и окружности имеют уравнения вида
		\[
			A(z\bar{z}) + (B+Ci)z + (B-Ci)\bar{z} + D = 0,
		\]
		где $A, B, C, D \in \mathbb{R}$.
		Достаточно показать, что при $z = \frac{a w + b}{cw + d}$ это уравнение переходит в уравнение того же вида.
 
	\end{enumerate}
 
	\subsection{3.2.2. Геометрия Лобачевского}
 
	\begin{enumerate}[resume*]
		\item Указание. Прямые и окружности имеют уравнения вида $A(z\bar{z}) + \overline{(B+Ci)z} + \ldots$
 
	\end{enumerate}
\end{document}
